\section{The Directives, and What Each One Is For}
\label{sec:ch05-the-directives}

\index{directives!the vocabulary|(}%
Federation v2 hands every subgraph library a block of sixteen directive
definitions to add to your schema for you, plus one more for libraries that
cannot write GraphQL's own \code{extend}
keyword~\autocite{apollosubgraphspec}. Four of the sixteen carry the model.
The rest get a paragraph each so that you recognize one in somebody else's
schema, and appendix~\ref{app:directives} carries the definitions in full.

\subsection{The Opt-In}

\index{link@\code{"@link}!is the version declaration}%
\code{@link} comes first because nothing else exists without it. It goes on the
schema itself and says which version of the vocabulary this document is written
in and which of its names the document uses:

\begin{graphqlsdl}
directive @link(url: String!, as: String, for: link__Purpose,
                import: [link__Import]) repeatable on SCHEMA
\end{graphqlsdl}

\index{link@\code{"@link}!the import list}%
The url ends in a version, and a subgraph that names an older one gets the
vocabulary of that older one, which is how a graph composes while its services
are on different versions. The import list does the other job: a directive you
did not import is not available under its bare name. In practice the library
writes this line for you from what your schema actually uses, and the first
time you will care is the day you reach for a directive your version does not
have and get an error about an unknown name rather than about a version.

\index{link@\code{"@link}!Hot Chocolate ships a narrower one}%
One caution. The definition above is the specification's. The one Hot
Chocolate's federation package writes into an exported schema has neither
\code{as} nor \code{for}, and types \code{import} as a plain list of strings, so
the namespacing \code{as} exists for is not available to you here. Nothing in
this book needs it, and I say so because reading the specification and
then reaching for that argument is a natural thing to do.

\subsection{The Seam}

\index{directives!the four that describe a seam}%
Four directives describe what crosses a boundary, and
section~\ref{sec:ch05-an-entity-is-a-type-with-a-key} has already spent itself
on the first. \code{@key} makes a type an entity. The other three qualify what
a subgraph may say about an entity it does not own:

\begin{graphqlsdl}
directive @external on FIELD_DEFINITION | OBJECT
directive @requires(fields: FieldSet!) on FIELD_DEFINITION
directive @provides(fields: FieldSet!) on FIELD_DEFINITION
\end{graphqlsdl}

\index{external@\code{"@external}!declares a field this subgraph does not own}%
\code{@external} marks a field that appears in this subgraph's schema but is
owned somewhere else. It is a declaration of borrowing: the field is here so
that something in this schema can refer to it, and this service does not
resolve it.

\index{requires@\code{"@requires}!asks the router for data first}%
\code{@requires} is what borrowing is usually for. A field you own can need a
value you do not have, and \code{@requires} names that value as a field set, so
the router fetches it from its owner and hands it to your resolver before your
resolver runs. It turns one fetch into two that must happen in order, and that
ordering is the price of the convenience.

\index{provides@\code{"@provides}!offers a field the router would otherwise fetch}%
\code{@provides} runs the other way. It says that along this particular field,
this subgraph can already supply some of the target entity's fields, so the
router may skip a hop it was going to make. It is an optimization with a
promise attached, and the promise is the same one \code{@shareable} makes: the
value has to match what the owner would have said.

Chapter~\ref{ch:second-third-subgraph} is where the first two get used on a
real seam and the third is weighed and left out, and
chapter~\ref{ch:satisfiability} is where getting them wrong produces
the composition errors that are hardest to read.

\subsection{Overlap and Visibility}

\index{directives!the four that manage overlap}%
\code{@shareable} and \code{@override} are
section~\ref{sec:ch05-who-owns-a-field}'s, and two more belong beside them
because they also change what the supergraph shows rather than what a subgraph
computes.

\index{inaccessible@\code{"@inaccessible}!hides a field from clients}%
\code{@inaccessible} takes a field out of the client's reach without taking it
out of the graph. Apollo is precise about the distinction and it is one to
hold on to: the field \enquote{is not omitted from the supergraph schema, so
the router still knows it exists (but clients can't include it in
operations)}~\autocite{apollodirectives}. Other subgraphs may still reference
it. It is the way to add a field to one service before the others are ready to
have it exist, and the way to take one off the public graph without taking it
out of the code.

\index{tag@\code{"@tag}!labels a field for tooling}%
\code{@tag} \enquote{applies arbitrary string metadata to a schema
location}~\autocite{apollodirectives} and means nothing by itself. What reads
the tags is whatever you point at the graph: a tool that builds a restricted
view of the supergraph for one audience, a policy check, a report. It is a
hook, and the value comes from what you hang on it.

\subsection{Access Control}

\index{directives!the three that carry authorization}%
\code{@authenticated}, \code{@requiresScopes} and \code{@policy} let a subgraph
state a requirement in its schema as well as in its resolvers, so that the
router can enforce it too. Chapter~\ref{ch:entities-authorization}
is about the reason to care, which is not the one people expect: the entity
route is a second door into every object in your graph, and a guard written on
the field a client normally uses does not stand in front of it.

\subsection{The Rest}

\index{interfaceObject@\code{"@interfaceObject}!contributes to every implementation}%
\code{@interfaceObject} lets a subgraph add a field to every type implementing
an interface without knowing what those types are. Apollo describes the
subgraph's declaration as \enquote{an abstraction of another subgraph's entity
interface}~\autocite{apollodirectives}: you declare an object type carrying the
interface's name and its key, and composition distributes the field to each
implementation.

\index{composeDirective@\code{"@composeDirective}!carries a custom directive through}%
\code{@composeDirective} names one of your own directives and asks the composer
to carry it into the supergraph. It is needed because \enquote{by default,
composition omits most directives from the supergraph
schema}~\autocite{apollodirectives}, which is the right default and an awkward
surprise the first time a directive of yours vanishes.

\index{extends@\code{"@extends}!for libraries without the keyword}%
\code{@extends} is the seventeenth, and it exists only for libraries that
cannot write GraphQL's own \code{extend} keyword. Hot Chocolate can, so this
book never uses it, and it is here because you will meet it in schemas written
by libraries that cannot.

\index{directives!what this package does not ship}%
Two of the sixteen are not in the package this book uses. \code{@context} and
\code{@fromContext}, which pass a value down from an ancestor field into a
resolver's argument, have no attribute and no directive type in
\code{HotChocolate.ApolloFederation}, so no chapter here claims anything about
them.

\index{cost@\code{"@cost}!arrives from elsewhere}%
\code{@cost} and \code{@listSize} are a different case again. Apollo documents
both alongside the sixteen, as controls on how much work a client may ask
for rather than as anything to do with a seam, and Hot Chocolate implements
them in its core rather than in the federation package. Which is why
chapter~\ref{ch:data-without-n-plus-1}'s exported schema was already carrying
both before this book had said the word subgraph.
\index{directives!the vocabulary|)}%
